\documentclass[12pt]{article}
\usepackage[T1]{fontenc}
\usepackage[utf8]{inputenc}
\usepackage{lmodern}
\usepackage{textcomp}
\usepackage{enumitem}
\usepackage{geometry}
\usepackage{graphicx}
\usepackage{amsmath, amssymb}
\geometry{margin=1in}
\begin{document}
\begin{center}
Constitutional Law - Undergraduate Level Assessment 18 \
Total Marks: 40 \
Answer all questions.
\end{center}

\section*{Multiple Choice (10 marks)}
\begin{enumerate}[label=\arabic*.]
\item What is the protective principle of jurisdiction?
\begin{enumerate}[label=(\alph*)]
    \item It is jurisdiction based on the harm to national interests by conduct committed abroad
    \item It is jurisdiction in order to protect one's nationals abroad
    \item It is jurisdiction in order to protect international human rights
    \item It is jurisdiction based on the nationality of the offender
\end{enumerate}
\item Critical Legal Studies (CLS) is often compared to (or with) American realism. Which of the following statements is inaccurate?
\begin{enumerate}[label=(\alph*)]
    \item Both are concerned with the 'law in action.'
    \item Both are anti-formalist and sceptical
    \item Both adopt a liberal ideology.
    \item Both attempt to demystify the law.
\end{enumerate}
\item Soft positivism accepts that the rule of recognition may include moral criteria.' Which proposition below is the most inconsistent with this description?
\begin{enumerate}[label=(\alph*)]
    \item Incorporationism' accepts that judges may decide cases by reference to moral factors.'
    \item A soft positivist rejects the role of morality in the description of law.
    \item Sometimes the definition of law includes moral considerations.
    \item Judges lack strong discretion.
\end{enumerate}
\item Which philosopher called the idea of natural rights 'nonsense on stilts'?
\begin{enumerate}[label=(\alph*)]
    \item Alan Gerwith
    \item Emmanuel Kant
    \item John Locke
    \item Jeremy Bentham
\end{enumerate}
\item Which of these statements best describes the UK Constitution?
\begin{enumerate}[label=(\alph*)]
    \item The UK Constitution's only source of power is that of the sovereign
    \item The UK Constitution gives the judiciary the power to overturn acts of parliament
    \item The UK Constitution is uncodified and can be found in a number of sources
    \item The UK Constitution is based on a Bill of Rights
\end{enumerate}
\end{enumerate}

\section*{True / False (10 marks)}
\textit{Answer True or False}
\begin{enumerate}[label=\arabic*.]
\setcounter{enumi}{5}
\item True or False: The Critical Legal Studies School argues that legal rules serve as tools for the powerful to perpetuate the status quo.
\item True or False: Fuller's concept of law aligns closely with a positivist view, emphasizing strict adherence to established legal rules.
\item True or False: Kant's 'categorical imperative' requires individuals to always act in the best interests of their community over personal values.
\item True or False: Recognition of governments is primarily granted only to democratic entities in contemporary international relations.
\item True or False: The rule of recognition confers power and imposes duties on judges to decide cases within a legal system.
\end{enumerate}

\section*{Long Form Response (20 marks)}
\textit{Answer each question with a clear and complete explanation.}
\begin{enumerate}[label=\arabic*.]
\setcounter{enumi}{10}
\item Discuss the common misconceptions surrounding Maine's aphorism 'the movement of progressive societies has hitherto been a movement from Status to Contract' and analyze its implications in constitutional law.
\item Discuss the relationship between customary law and treaty provisions, analyzing how treaties can influence the formation, codification, and evolution of customary law.
\end{enumerate}

\vfill
\noindent\textit{End of Paper}
\end{document}